\section{Introduction}
\label{sec:introduction}

Mobile agents that answer questions about their surroundings must do three things well: \emph{explore} efficiently, \emph{remember} what they have seen in a form useful for reasoning, and \emph{update} that memory when the world changes.
Embodied question answering (EQA) makes this explicit: the agent receives a question (e.g., ``What color is the mug on the counter?'') and must explore until it can answer with calibrated confidence.

Two recent lines of work address complementary parts of this problem.
\textbf{GraphEQA}~\citep{saxena2024grapheqa} maintains a real-time 3D semantic scene graph and task-relevant images as multi-modal memory for a vision-language model (VLM), with hierarchical exploration in Habitat on HM-EQA and OpenEQA-style benchmarks.
\textbf{DynaMem}~\citep{liu2024dynamem} maintains a dynamic voxel semantic map that supports adding and removing geometry as objects move, enabling open-vocabulary mobile manipulation in changing homes.

\textbf{Dynagraph} unifies these ideas for \emph{exploring robot agents}: the same observations feed a voxel map used for navigation (DynaMem) and an object graph used for EQA (GraphEQA), with explicit \textbf{graph lifecycle} rules---spatial merge of duplicate detections and staleness-based pruning---so the graph does not grow without bound and does not retain objects that are no longer supported by perception.

\subsection{Contributions}

\begin{itemize}
  \item \textbf{Dynagraph memory architecture}: coupled voxel navigation map + dynamic object scene graph with merge and staleness (\texttt{dynagraph\_merge\_xy\_m}, \texttt{dynagraph\_staleness\_horizon}).
  \item \textbf{Open-source implementation} in \texttt{emet} (\texttt{emet run dynagraph}), building on \texttt{GraphEQAController} and \texttt{SparseVoxelMap}.
  \item \textbf{Habitat EQA harness} (separate development branch) to reproduce GraphEQA-style evaluation on HM-EQA / OpenEQA.
  \item \textbf{Extended MuJoCo evaluation} on Robocasa and MolmoSpaces across Stretch, Galaxea R1, and optional registry robots.
\end{itemize}

\subsection{Limitations (draft)}

Habitat harness and full experimental results are in progress.
Our graph builder differs from Hydra-based 3DSG pipelines used in the original GraphEQA stack; we compare on EQA outcomes, not intermediate graph topology.
Real-home manipulation trials at OK-Robot or DREAM scale are left to future work.
